% -----------------------------------------------------------------------------------------------------
\section{双部分子散射}\label{sec:front-dps}

一次 $pp$ 碰撞中的两个硬散射（如双 Drell--Yan）对积分截面形式上是幂次压低的，但对横动量微分截面（$q_{T,i}\ll Q_i$）是领头的，在小 $x$ 多末态时还会增强。框架（据各文摘要）：色关联的双部分子分布（DPD）、幂次计数、全阶证明的要素以及与 GPD 和格点的联系\cite{DiehlOstermeierSchafer2012}（连同勘误）；色关联、干涉、演化、$1\to2$ 劈裂奇异性（小间距 $y$ 处 $\propto1/y^2$）与重复计数问题\cite{Manohar:2012jr,Manohar:2012pe}。双 Drell--Yan 中积分截面与 $q_T$ 微分截面的全阶 Glauber 相消\cite{Diehl:2015bca}，以及借非阿贝尔 Ward 恒等式使软胶子退耦为 Wilson 线真空期望值\cite{Diehl:2018wfy}，属于核心陈述\cref{thm:break-double-dy}（[A]）。

\begin{proposition}[无重复计数的双部分子散射 \tagA]\label{prop:front-dps-scheme}
\begin{equation}\label{eq:front-dgs}
\sigma=\sigma_{\rm SPS}+\sigma_{\rm DPS}-\sigma_{\rm sub},\qquad
\frac{\dd\sigma_{\rm DPS}}{\dd x_1\dd x_2\dd\bar x_1\dd\bar x_2}=\sum\hat\sigma_{a_1b_1}\hat\sigma_{a_2b_2}\!\int\!\dd^2y\,\Phi^2(y\nu)F_{b_1b_2}(\bar x_i,y)F_{a_1a_2}(x_i,y),
\end{equation}
其中截断函数满足 $u\gg1$ 时 $\Phi(u)\to1$、小 $u$ 时 $\Phi(u)=\order{u^\delta}$（如 $\Theta(u-b_0)$ 或 $1-\ee^{-u^2/4}$），$\nu\sim Q$；$\sigma_{\rm sub}$ 是把两个 DPD 都换成其固定阶微扰劈裂形式的 DPS 公式（对 1v1 即 $\sigma_{1v1,{\rm pt}}$）。小 $y$ 处 $\Phi$ 依赖在 $\sigma_{\rm DPS}$ 与 $\sigma_{\rm sub}$ 之间相消（剩下 SPS），大 $y$ 处 $\sigma_{\rm SPS}\simeq\sigma_{\rm sub}$（剩下 DPS）。由此 DPD 被定义为可在高阶使用的算符矩阵元\cite[第~4.1~节]{Diehl:2017kgu}。非微扰地位：\cref{asm:intro-composite}、\cref{asm:intro-bridge}（MA0、MA1+MA2）。
% VF-FRO-13
% errtype: E-S
\uses{def:break-dpdf, thm:break-double-dy}
\end{proposition}

\begin{proofsketch}
把 DPS 区域与 SPS 中 $1\to2$ 劈裂区域的重叠写成以 $\Phi$ 截断的 $y$ 积分，在小 $y$ 处用 DPD 的微扰劈裂形式构造减除项 $\sigma_{\rm sub}$，逐阶检验两端极限中 $\Phi$ 依赖的相消。完整论证见\cite[第~4~节]{Diehl:2017kgu}。
\end{proofsketch}

\paragraph{地位与缺口。}[A]：方案的自洽性在固定阶论证；DPD 演化之后，相消只在更高阶意义下成立。带色末态或喷注的完整 DPS 因子化仍是猜想（\cref{conj:break-dps-coloured}）。

\begin{proposition}[多部分子软因子快度发散的重整化定理 \tagA]\label{prop:front-rapidity-renorm}
多部分子软因子的快度发散来自与无穷远横平面的交换，经一个共形映射在结构上等价于紫外发散；由此得到快度发散的重整化定理（借共形变换证明）、多部分子 TMD 的定义，以及软反常维度与快度反常维度之间的精确关系（并导出三圈快度反常维度）（\cite{Vladimirov:2017ksc}，据该文摘要）。（本命题只涉及软因子即 Wilson 线真空期望值。）
\uses{def:tmd-soft, def:break-dpdf}
\end{proposition}

\begin{proofsketch}
用共形变换把类光 Wilson 线软因子中远横平面处的快度发散映到紫外发散，再搬用 Wilson 线算符的紫外重整化定理。完整论证见\cite{Vladimirov:2017ksc}。
\end{proofsketch}

\paragraph{地位与缺口。}[A]（只读到摘要层次）。第二十二、二十三章关于快度发散线性性的讨论指向这里；Vladimirov 称 SIDIS 软因子的快度因子化是“the open question”（据\cite{Liu:2023pjb}引用的\cite{Vladimirov:2017ksc}期刊版第~32~页）。

\begin{proposition}[色关联 DPD 的演化与匹配 \tagP]\label{prop:front-dpd-evolution}
到次领头阶（两圈核），均据摘要：TMD DPS 的因子化、演化与匹配\cite{Buffing:2017mqm}（连同 2021 年的勘误／增补）；色非单态 DPD 的 NLO 演化核\cite{Diehl:2022rxb}；两圈 $1\to2$ 劈裂\cite{Diehl:2019rdh}；劈裂中的夸克质量效应\cite{Diehl:2022dia}；坐标空间中的动量与数目求和规则\cite{Diehl:2020xyg}；综述\cite{Diehl:2017wew}。（本命题只涉及演化核与部分子层次的匹配系数。）
\uses{def:break-dpdf, prop:front-rapidity-renorm}
\end{proposition}

\paragraph{地位与缺口。}[P]（算到 NLO）；色关联 DPD 的唯象仍是开放问题。
